\chapter{Expressions and definitions}
\label{ch:expressions}

\section{Programs}
\index{program}
A program is a sequence of top-level items: definitions (\fc|let|,
\fc|let rec|), type declarations (\fc|type|), traits and instances
(\fc|trait|, \fc|impl|), modules (\fc|mod|, \fc|use|), imports and macros.
Items are processed in order, and each can use what the items before it
define. Expressions appear inside definitions; to run an expression for its
effect, bind it to the unit pattern: \fc|let () = ...|.

\section{Definitions}
\label{sec:let}
\index{let@\texttt{let}}\index{let rec@\texttt{let rec}}

\fc|let x = e| binds the value of \fc|e| to \fc|x|. A definition with
parameters, \fc|let f x y = e|, is shorthand for
\fc|let f = fn x -> fn y -> e|. The left-hand side can be any pattern
(Chapter~\ref{ch:patterns}): \fc|let (q, r) = (7 / 2, 7 % 2)|.
\fc|let rec| makes the name visible in its own definition, and \fc|and|
joins mutually recursive definitions.

Inside an expression, \fc|let p = e1 in e2| binds \fc|p| for \fc|e2| only.
Definitions are \emph{polymorphic}: a name bound by \fc|let| can be used at
several types (Section~\ref{sec:inference}).

\forgeexample{expr-let}

\section{Functions}
\label{sec:functions}
\index{functions}

\fc|fn p -> e| is a function of one argument; \fc|fn x -> fn y -> e| a
curried function of two. Application is juxtaposition, \fc|f x y|, and
associates to the left. A parameter is a pattern, and parameters and results
may carry type annotations: \fc|fn (x : Int) : Int -> x + 1|, or in a
definition \fc|let scale (k : Int) (x : Int) : Int = k * x|. Functions are
closures: they capture the variables in scope where they are created.

Every binary operator is also a function, written in parentheses:
\fc|(+)|, \fc|(::)|, \fc|(^)|; \fc|(*) 10| is the function that multiplies
by ten.

\forgeexample{expr-functions}

\section{Conditionals and logic}
\index{if@\texttt{if}}

\fc|if c then a else b| requires \fc|c : Bool| and gives \fc|a| and \fc|b|
the same type; there is no \fc|if| without \fc|else|. \fc|&&| and \fc!||!
evaluate their right operand only when needed, and \fc|not| negates.

\section{Pattern matching with \texttt{case}}
\index{case@\texttt{case}}
\fc@case e do |@ \textit{p}$_1$ \fc@->@ \textit{e}$_1$ \fc@| ...@ evaluates
\fc|e| and the first arm whose pattern matches. If no arm matches, the
program stops with the run-time error \textit{no pattern matched}; the
compiler warns when that can happen (Section~\ref{sec:exhaustiveness}).

\forgeexample{expr-case}

\section{Operators}
\label{sec:operators}
\index{operators!precedence}

Table~\ref{tab:precedence} lists the operators from the loosest to the
tightest binding. A custom operator takes the precedence and associativity
of the fixed operators that share its first character.

\begin{table}[ht]
\centering\small
\caption{Operator precedence, loosest first.}
\label{tab:precedence}
\begin{tabularx}{\textwidth}{@{}l l X@{}}
\toprule
\textbf{Operators} & \textbf{Associativity} & \textbf{Meaning} \\
\midrule
\tc{|>} & left & pipeline: \tc{x |> f} is \tc{f x} \\
\tc{::} & right & cons \\
\tc{||} & right & or (short-circuit) \\
\texttt{\&\&} & right & and (short-circuit) \\
\tc{== != <> < > <= >=} & left & comparison; custom operators starting with \tc{<}, \tc{>}, \tc{=} or \tc{!} \\
\tc{+ - ^} & left & addition, subtraction, string concatenation; custom operators starting with \tc{+}, \tc{-} or \tc{^} \\
\tc{* / \%} & left & multiplication, division, remainder; custom operators starting with \tc{*}, \tc{/} or \tc{\%} \\
application \tc{f x} & left & \\
\tc{-e}, \tc{e.field} & & negation, field access \\
\bottomrule
\end{tabularx}
\end{table}

\paragraph{Arithmetic} \fc|+ - * / %| work on \fc|Int|, and \fc|+. -. *. /.| on \fc|Float| (with the precedence of \fc|+| and \fc|*|); integer division truncates
toward zero, \fc|%| has the sign of its left operand, and dividing by zero is
a run-time error.\label{sec:int-width}\index{integers!width}
Integers wrap around on overflow. Their width is not yet the same in the two
back ends: 63 bits in the interpreter, 64 bits in native code, so results
that overflow 63 bits differ (an open question of the implementation).

\paragraph{Comparison} \fc|==|, \fc|!=| (also written \fc|<>|), \fc|<|,
\fc|>|, \fc|<=| and \fc|>=| are \emph{structural} and work on values of any
type except functions: numbers and characters by value, strings
lexicographically, lists and tuples element by element, records field by
field, and values of a sum type first by constructor, in declaration order,
then by payload. Comparing functions is a run-time error.

\paragraph{Strings} \fc|^| concatenates two strings.

\forgeexample{expr-operators}

\section{Lists}
\index{lists}\index{comprehensions}
A list is written \fc|[1, 2, 3]| or built with \fc|::|; \fc|[]| is the
empty list. \fc|[a...b]| is the list of the integers from \fc|a| to \fc|b|
(empty if \fc|a > b|). A \emph{comprehension}
\fc@[e | p <- xs, cond, q <- ys]@ evaluates \fc|e| for every combination of
generators, from left to right; a generator's pattern may be refutable, and
the elements it does not match are skipped; a condition keeps the
combinations for which it holds.

\forgeexample{expr-lists}

\section{Tuples and records}
\index{tuples}\index{records}
A tuple groups values of possibly different types: \fc|(1, "a", true)|. A
record names its fields: \fc|{ x: 1, y: 2 }|. \fc|r.x| reads a field and
\fc|{ r with x = 5 }| is a copy of \fc|r| with a new value for \fc|x|. Record
types are structural (Section~\ref{sec:records}); a record type can be
named with \fc|type|, which is optional.

\forgeexample{expr-records}

\section{Blocks}
\index{blocks}
A block \fc|{ s1; s2; ...; e }| is a sequence of local definitions and
statements, separated by semicolons. A definition scopes over the rest of
the block, a statement is evaluated for its effect, and the block's value is
that of its last part (\fc|()| if that is a definition).

\forgeexample{expr-blocks}

\section{Pipelines}
\index{pipeline}
\fc!x |> f! applies \fc|f| to \fc|x|. Chains of transformations read from
left to right:

\forgeexample{expr-pipeline}
